% =====================================================================
%  method.tex  --  PHASE VII, Task 14C
%
%  Section 3.  The framework overview, the data, and the two forecasting
%  formulations.  The state-inference subsection lives in its own file
%  and is pulled in here.
%
%  Intended placement: Section 3 (\label{sec:method}), after related
%  work.
%
%  Drop-in usage: % =====================================================================
%  method.tex  --  PHASE VII, Task 14C
%
%  Section 3.  The framework overview, the data, and the two forecasting
%  formulations.  The state-inference subsection lives in its own file
%  and is pulled in here.
%
%  Intended placement: Section 3 (\label{sec:method}), after related
%  work.
%
%  Drop-in usage: % =====================================================================
%  method.tex  --  PHASE VII, Task 14C
%
%  Section 3.  The framework overview, the data, and the two forecasting
%  formulations.  The state-inference subsection lives in its own file
%  and is pulled in here.
%
%  Intended placement: Section 3 (\label{sec:method}), after related
%  work.
%
%  Drop-in usage: % =====================================================================
%  method.tex  --  PHASE VII, Task 14C
%
%  Section 3.  The framework overview, the data, and the two forecasting
%  formulations.  The state-inference subsection lives in its own file
%  and is pulled in here.
%
%  Intended placement: Section 3 (\label{sec:method}), after related
%  work.
%
%  Drop-in usage: \input{method}
%  This file \inputs method_gmm_labeller.tex.
%  Bibliography: paper/references.bib
% =====================================================================

\section{Proposed Framework}
\label{sec:method}

\subsection{Overview}
\label{sec:method:overview}

The framework is a three-stage pipeline over a single input: the
per-phase voltage and current a standard retrofit meter already
records. Each stage is specified by what the next one needs.

\begin{enumerate}
\item \emph{State inference} (Section~\ref{sec:method:labelling}) maps
  the electrical record to a label per sample from
  $\{$OFF, STANDBY, WORKING, PEAK\_LOAD$\}$ without annotations, and is
  required to produce a sequence whose state durations are physically
  admissible --- not merely a per-sample partition --- because duration
  forecasting over a flickering label sequence is ill-posed.
\item \emph{Duration forecasting} (Sections~\ref{sec:forecasting}
  and~\ref{sec:decision:policies}) predicts how much longer the machine
  will remain non-productive, and is required to target the conditional
  \emph{mean} of that duration, because the decision rule below is
  affine in it.
\item \emph{Decision} (Section~\ref{sec:decision}) converts the forecast
  into a de-energise/keep-running action by minimising total cost
  subject to thermal and wear constraints, and returns the break-even
  duration as a derived quantity rather than consuming it as a
  parameter.
\end{enumerate}

The stages are separable in implementation and are evaluated separately
throughout, which is what makes the component ablation of
Section~\ref{sec:ablation} meaningful.

\subsection{Data}
\label{sec:method:data}

IMDELD~\cite{imdeld2018dataset,martins2018imdeld} publishes $1$\,Hz
three-phase measurements from eight machines of a Brazilian
poultry-feed pellet plant: two pelletizers, two milling machines, two
exhaust fans and two double-pole contactors. Six channels are used ---
active, reactive and apparent power, current, voltage and power factor
--- and no other information about the machines enters the models. The
plant's published calendar (weekday operation, a halt between 17:00 and
22:00 local time to avoid the peak tariff, and no weekend production)
is used \emph{only} as external validation evidence, never as a model
input to the state layer.

Coverage differs sharply by machine and is reported as measured rather
than as published. The pelletizer~I record spans $155$ days from
2017-10-30 to 2018-04-03 and contains $63.4$ days of data in $74$
contiguous segments; $31$ acquisition gaps account for the remainder,
the largest of them $30$ days. The two milling machines cover $12.2$ of
$43$ days in $1\,082$ and $580$ fragments. \textbf{Every annualised
quantity in this paper is scaled by covered time, not by span}, which
matters by a factor of $2.4$ on the five-month records and is the
reason the milling machines' figures are not annualised at all
(Section~\ref{sec:limitations:milling}).

\input{method_gmm_labeller}

\subsection{Two forecasting problems, not one}
\label{sec:method:forecasting}

The framework contains two forecasters. They are frequently conflated,
they consume different features, and no result in this paper reports
them in the same column.

\emph{The window state forecaster} takes a $600$\,s lookback of the
electrical record and predicts the state at each step of a $300$\,s
horizon. It is a sequence problem, it is what the baseline comparison of
Section~\ref{sec:forecasting} evaluates, and it is scored on windows.
Its design representation summarises each channel over the lookback ---
final observed value, sub-block means and standard deviations --- giving
$217$ columns for the tree, or the raw sequence for the neural
architectures.

\emph{The episode-duration forecaster} takes the state of an idle
episode at a decision epoch and predicts the \emph{remaining} idle time
$R$ in seconds. It is a regression problem, it is what the decision
layer of Section~\ref{sec:decision} consumes, and it is scored per
decision epoch. Its ten features are elapsed idle time, the current
position in the factory calendar, and the load history of the
production run that just ended.

Keeping them apart is not pedantry. Section~\ref{sec:explain} shows they
draw on almost disjoint information --- $11.6\%$ of the window model's
attribution is calendar against $82.4\%$ of the duration model's --- and
Section~\ref{sec:crossmachine} shows that one transfers between machines
and the other does not, for exactly that reason. Attributing the
economic result to the window model's $F_1$ would credit a model that
plays no part in producing it.

%  This file \inputs method_gmm_labeller.tex.
%  Bibliography: paper/references.bib
% =====================================================================

\section{Proposed Framework}
\label{sec:method}

\subsection{Overview}
\label{sec:method:overview}

The framework is a three-stage pipeline over a single input: the
per-phase voltage and current a standard retrofit meter already
records. Each stage is specified by what the next one needs.

\begin{enumerate}
\item \emph{State inference} (Section~\ref{sec:method:labelling}) maps
  the electrical record to a label per sample from
  $\{$OFF, STANDBY, WORKING, PEAK\_LOAD$\}$ without annotations, and is
  required to produce a sequence whose state durations are physically
  admissible --- not merely a per-sample partition --- because duration
  forecasting over a flickering label sequence is ill-posed.
\item \emph{Duration forecasting} (Sections~\ref{sec:forecasting}
  and~\ref{sec:decision:policies}) predicts how much longer the machine
  will remain non-productive, and is required to target the conditional
  \emph{mean} of that duration, because the decision rule below is
  affine in it.
\item \emph{Decision} (Section~\ref{sec:decision}) converts the forecast
  into a de-energise/keep-running action by minimising total cost
  subject to thermal and wear constraints, and returns the break-even
  duration as a derived quantity rather than consuming it as a
  parameter.
\end{enumerate}

The stages are separable in implementation and are evaluated separately
throughout, which is what makes the component ablation of
Section~\ref{sec:ablation} meaningful.

\subsection{Data}
\label{sec:method:data}

IMDELD~\cite{imdeld2018dataset,martins2018imdeld} publishes $1$\,Hz
three-phase measurements from eight machines of a Brazilian
poultry-feed pellet plant: two pelletizers, two milling machines, two
exhaust fans and two double-pole contactors. Six channels are used ---
active, reactive and apparent power, current, voltage and power factor
--- and no other information about the machines enters the models. The
plant's published calendar (weekday operation, a halt between 17:00 and
22:00 local time to avoid the peak tariff, and no weekend production)
is used \emph{only} as external validation evidence, never as a model
input to the state layer.

Coverage differs sharply by machine and is reported as measured rather
than as published. The pelletizer~I record spans $155$ days from
2017-10-30 to 2018-04-03 and contains $63.4$ days of data in $74$
contiguous segments; $31$ acquisition gaps account for the remainder,
the largest of them $30$ days. The two milling machines cover $12.2$ of
$43$ days in $1\,082$ and $580$ fragments. \textbf{Every annualised
quantity in this paper is scaled by covered time, not by span}, which
matters by a factor of $2.4$ on the five-month records and is the
reason the milling machines' figures are not annualised at all
(Section~\ref{sec:limitations:milling}).

% =====================================================================
%  method_gmm_labeller.tex  --  PHASE V CORRECTION 4 & 5
%
%  Drop-in description of the GMM-HMM labelling procedure, including
%  the accept/reject step that Phase V Task 12B showed is load-bearing
%  for reproducibility.  Insert into Section 3 (Method).
%
%  Numbers: all STANDBY-time totals are rounded to 2 significant
%  figures and stated as conditional on the physics battery passing,
%  per Phase V Task 11C (block-bootstrap CI ±15–25%) and
%  Task 12B (labelling CV 33.7% unconditional, 3.0% conditional).
% =====================================================================

\subsection{Machine-state inference}
\label{sec:method:labelling}

\subsubsection{Preprocessing}
\label{sec:method:preprocess}

Each machine's $1$\,Hz electrical record is preprocessed in three
deterministic steps, applied identically to every machine so that
differences in downstream results are attributable to the machines,
not to the preprocessing.

\begin{enumerate}
\item \emph{Spike masking.}  For each channel, rolling-MAD outliers
  (window $11$, threshold $5\times$MAD) are flagged and excluded
  from model fitting but retained in the record; removing them would
  destroy the row-count-to-seconds correspondence that episode
  durations depend on.

\item \emph{OFF-state denoising.}  Rows with active power below
  $5$\,W are smoothed with a rolling median (window $5$) to suppress
  sensor noise in the de-energised state without touching on-state
  dynamics.

\item \emph{Feature standardisation.}  Six electrical channels are
  log$_1$p-transformed (right-skewed channels: $P, Q, S, I$) and
  z-scored jointly before any model fitting.
\end{enumerate}

\subsubsection{Gaussian hidden Markov model}
\label{sec:method:ghmm}

Emission distributions are estimated by a Gaussian mixture model
(GMM, full covariance, $k = 4$ components) fitted on a uniform
random sample of $200{,}000$ rows.  A Gaussian HMM is then
warm-started from these emissions with $\texttt{params}=\texttt{``t''}$,
so only the transition matrix is learned from $64$ contiguous
two-hour blocks spread across the record; this decouples emission
shape (a property of the marginal distribution) from transition
dynamics (a property of the sequence) and avoids the local optima
that Baum--Welch finds when initialised randomly on this
dataset~\cite{fox2011sticky}.  Viterbi decoding is applied per
contiguous segment to prevent the decoder from asserting a transition
across an acquisition gap.  A minimum-dwell constraint of $10$\,s
is applied post-decode to remove label flicker without reshaping
the partition (Section~\ref{sec:results:task4}).

\subsubsection{Physics-validation accept/reject loop}
\label{sec:method:accept_reject}

\emph{This step is load-bearing for reproducibility and was not
explicit in the Phase I--IV description; it is stated here as a
result of Phase~V Task~12B.}

After decoding, a battery of ten physics tests (T1--T10:
power-factor separation, no-load current ratio, power conservation,
dwell plausibility, cross-machine replication, and factory-schedule
agreement) is run on a strided sample of the labelled record and
aggregated into a physics score $s \in [0, 1]$.  If $s < 0.90$,
the model is re-fitted with a different seed (incremented by a
fixed prime) and the battery is re-evaluated, up to three attempts.

The motivation is empirical.  Phase~V Task~12B fitted the labeller at
ten different seeds on identical preprocessed data from pelletizer~I
and found that three seeds (30\%) converge to a wrong cluster:

\begin{itemize}
\item seed~6: the ``STANDBY'' cluster absorbed a loaded running state
  (mean power $12{,}961$\,W instead of the correct $\sim3{,}750$\,W);
\item seeds~8 and~9: the ``STANDBY'' cluster absorbed the OFF state
  (mean power $0.8$--$1.1$\,W).
\end{itemize}

The physics battery scored those three seeds at $0.60$--$0.80$ and the
seven correct seeds at $1.00$, making the accept/reject decision exact
on this record.  Conditioning on a passing seed collapses the
STANDBY-hour coefficient of variation from $33.7\%$ (all ten seeds)
to $3.0\%$ (passing seeds), and the total reported throughout this
paper---$92 \pm 3$\,h over the 155-day record---is the mean of the
seven passing seeds.

\emph{All STANDBY-time totals in this paper are therefore stated as
conditional on the physics battery passing} (i.e., $s \ge 0.90$)
\emph{and are rounded to two significant figures}, because the
block-bootstrap sampling uncertainty (Phase~V Task~11C, one-week
blocks) is $\pm 15$--$25\%$ and the conditional labelling
uncertainty is $\pm 3\%$.  Three significant figures are not
supported by the data.

The accept/reject loop is implemented in \texttt{src/labelling.py}
(\texttt{label\_record}, \texttt{physics\_score\_threshold}=0.90,
\texttt{max\_refit\_attempts}=3) and its output is stored in the
model-info dictionary (\texttt{physics\_validation\_attempts},
\texttt{physics\_validation\_seed}, \texttt{physics\_score}) for
audit and reproduction.


\subsection{Two forecasting problems, not one}
\label{sec:method:forecasting}

The framework contains two forecasters. They are frequently conflated,
they consume different features, and no result in this paper reports
them in the same column.

\emph{The window state forecaster} takes a $600$\,s lookback of the
electrical record and predicts the state at each step of a $300$\,s
horizon. It is a sequence problem, it is what the baseline comparison of
Section~\ref{sec:forecasting} evaluates, and it is scored on windows.
Its design representation summarises each channel over the lookback ---
final observed value, sub-block means and standard deviations --- giving
$217$ columns for the tree, or the raw sequence for the neural
architectures.

\emph{The episode-duration forecaster} takes the state of an idle
episode at a decision epoch and predicts the \emph{remaining} idle time
$R$ in seconds. It is a regression problem, it is what the decision
layer of Section~\ref{sec:decision} consumes, and it is scored per
decision epoch. Its ten features are elapsed idle time, the current
position in the factory calendar, and the load history of the
production run that just ended.

Keeping them apart is not pedantry. Section~\ref{sec:explain} shows they
draw on almost disjoint information --- $11.6\%$ of the window model's
attribution is calendar against $82.4\%$ of the duration model's --- and
Section~\ref{sec:crossmachine} shows that one transfers between machines
and the other does not, for exactly that reason. Attributing the
economic result to the window model's $F_1$ would credit a model that
plays no part in producing it.

%  This file \inputs method_gmm_labeller.tex.
%  Bibliography: paper/references.bib
% =====================================================================

\section{Proposed Framework}
\label{sec:method}

\subsection{Overview}
\label{sec:method:overview}

The framework is a three-stage pipeline over a single input: the
per-phase voltage and current a standard retrofit meter already
records. Each stage is specified by what the next one needs.

\begin{enumerate}
\item \emph{State inference} (Section~\ref{sec:method:labelling}) maps
  the electrical record to a label per sample from
  $\{$OFF, STANDBY, WORKING, PEAK\_LOAD$\}$ without annotations, and is
  required to produce a sequence whose state durations are physically
  admissible --- not merely a per-sample partition --- because duration
  forecasting over a flickering label sequence is ill-posed.
\item \emph{Duration forecasting} (Sections~\ref{sec:forecasting}
  and~\ref{sec:decision:policies}) predicts how much longer the machine
  will remain non-productive, and is required to target the conditional
  \emph{mean} of that duration, because the decision rule below is
  affine in it.
\item \emph{Decision} (Section~\ref{sec:decision}) converts the forecast
  into a de-energise/keep-running action by minimising total cost
  subject to thermal and wear constraints, and returns the break-even
  duration as a derived quantity rather than consuming it as a
  parameter.
\end{enumerate}

The stages are separable in implementation and are evaluated separately
throughout, which is what makes the component ablation of
Section~\ref{sec:ablation} meaningful.

\subsection{Data}
\label{sec:method:data}

IMDELD~\cite{imdeld2018dataset,martins2018imdeld} publishes $1$\,Hz
three-phase measurements from eight machines of a Brazilian
poultry-feed pellet plant: two pelletizers, two milling machines, two
exhaust fans and two double-pole contactors. Six channels are used ---
active, reactive and apparent power, current, voltage and power factor
--- and no other information about the machines enters the models. The
plant's published calendar (weekday operation, a halt between 17:00 and
22:00 local time to avoid the peak tariff, and no weekend production)
is used \emph{only} as external validation evidence, never as a model
input to the state layer.

Coverage differs sharply by machine and is reported as measured rather
than as published. The pelletizer~I record spans $155$ days from
2017-10-30 to 2018-04-03 and contains $63.4$ days of data in $74$
contiguous segments; $31$ acquisition gaps account for the remainder,
the largest of them $30$ days. The two milling machines cover $12.2$ of
$43$ days in $1\,082$ and $580$ fragments. \textbf{Every annualised
quantity in this paper is scaled by covered time, not by span}, which
matters by a factor of $2.4$ on the five-month records and is the
reason the milling machines' figures are not annualised at all
(Section~\ref{sec:limitations:milling}).

% =====================================================================
%  method_gmm_labeller.tex  --  PHASE V CORRECTION 4 & 5
%
%  Drop-in description of the GMM-HMM labelling procedure, including
%  the accept/reject step that Phase V Task 12B showed is load-bearing
%  for reproducibility.  Insert into Section 3 (Method).
%
%  Numbers: all STANDBY-time totals are rounded to 2 significant
%  figures and stated as conditional on the physics battery passing,
%  per Phase V Task 11C (block-bootstrap CI ±15–25%) and
%  Task 12B (labelling CV 33.7% unconditional, 3.0% conditional).
% =====================================================================

\subsection{Machine-state inference}
\label{sec:method:labelling}

\subsubsection{Preprocessing}
\label{sec:method:preprocess}

Each machine's $1$\,Hz electrical record is preprocessed in three
deterministic steps, applied identically to every machine so that
differences in downstream results are attributable to the machines,
not to the preprocessing.

\begin{enumerate}
\item \emph{Spike masking.}  For each channel, rolling-MAD outliers
  (window $11$, threshold $5\times$MAD) are flagged and excluded
  from model fitting but retained in the record; removing them would
  destroy the row-count-to-seconds correspondence that episode
  durations depend on.

\item \emph{OFF-state denoising.}  Rows with active power below
  $5$\,W are smoothed with a rolling median (window $5$) to suppress
  sensor noise in the de-energised state without touching on-state
  dynamics.

\item \emph{Feature standardisation.}  Six electrical channels are
  log$_1$p-transformed (right-skewed channels: $P, Q, S, I$) and
  z-scored jointly before any model fitting.
\end{enumerate}

\subsubsection{Gaussian hidden Markov model}
\label{sec:method:ghmm}

Emission distributions are estimated by a Gaussian mixture model
(GMM, full covariance, $k = 4$ components) fitted on a uniform
random sample of $200{,}000$ rows.  A Gaussian HMM is then
warm-started from these emissions with $\texttt{params}=\texttt{``t''}$,
so only the transition matrix is learned from $64$ contiguous
two-hour blocks spread across the record; this decouples emission
shape (a property of the marginal distribution) from transition
dynamics (a property of the sequence) and avoids the local optima
that Baum--Welch finds when initialised randomly on this
dataset~\cite{fox2011sticky}.  Viterbi decoding is applied per
contiguous segment to prevent the decoder from asserting a transition
across an acquisition gap.  A minimum-dwell constraint of $10$\,s
is applied post-decode to remove label flicker without reshaping
the partition (Section~\ref{sec:results:task4}).

\subsubsection{Physics-validation accept/reject loop}
\label{sec:method:accept_reject}

\emph{This step is load-bearing for reproducibility and was not
explicit in the Phase I--IV description; it is stated here as a
result of Phase~V Task~12B.}

After decoding, a battery of ten physics tests (T1--T10:
power-factor separation, no-load current ratio, power conservation,
dwell plausibility, cross-machine replication, and factory-schedule
agreement) is run on a strided sample of the labelled record and
aggregated into a physics score $s \in [0, 1]$.  If $s < 0.90$,
the model is re-fitted with a different seed (incremented by a
fixed prime) and the battery is re-evaluated, up to three attempts.

The motivation is empirical.  Phase~V Task~12B fitted the labeller at
ten different seeds on identical preprocessed data from pelletizer~I
and found that three seeds (30\%) converge to a wrong cluster:

\begin{itemize}
\item seed~6: the ``STANDBY'' cluster absorbed a loaded running state
  (mean power $12{,}961$\,W instead of the correct $\sim3{,}750$\,W);
\item seeds~8 and~9: the ``STANDBY'' cluster absorbed the OFF state
  (mean power $0.8$--$1.1$\,W).
\end{itemize}

The physics battery scored those three seeds at $0.60$--$0.80$ and the
seven correct seeds at $1.00$, making the accept/reject decision exact
on this record.  Conditioning on a passing seed collapses the
STANDBY-hour coefficient of variation from $33.7\%$ (all ten seeds)
to $3.0\%$ (passing seeds), and the total reported throughout this
paper---$92 \pm 3$\,h over the 155-day record---is the mean of the
seven passing seeds.

\emph{All STANDBY-time totals in this paper are therefore stated as
conditional on the physics battery passing} (i.e., $s \ge 0.90$)
\emph{and are rounded to two significant figures}, because the
block-bootstrap sampling uncertainty (Phase~V Task~11C, one-week
blocks) is $\pm 15$--$25\%$ and the conditional labelling
uncertainty is $\pm 3\%$.  Three significant figures are not
supported by the data.

The accept/reject loop is implemented in \texttt{src/labelling.py}
(\texttt{label\_record}, \texttt{physics\_score\_threshold}=0.90,
\texttt{max\_refit\_attempts}=3) and its output is stored in the
model-info dictionary (\texttt{physics\_validation\_attempts},
\texttt{physics\_validation\_seed}, \texttt{physics\_score}) for
audit and reproduction.


\subsection{Two forecasting problems, not one}
\label{sec:method:forecasting}

The framework contains two forecasters. They are frequently conflated,
they consume different features, and no result in this paper reports
them in the same column.

\emph{The window state forecaster} takes a $600$\,s lookback of the
electrical record and predicts the state at each step of a $300$\,s
horizon. It is a sequence problem, it is what the baseline comparison of
Section~\ref{sec:forecasting} evaluates, and it is scored on windows.
Its design representation summarises each channel over the lookback ---
final observed value, sub-block means and standard deviations --- giving
$217$ columns for the tree, or the raw sequence for the neural
architectures.

\emph{The episode-duration forecaster} takes the state of an idle
episode at a decision epoch and predicts the \emph{remaining} idle time
$R$ in seconds. It is a regression problem, it is what the decision
layer of Section~\ref{sec:decision} consumes, and it is scored per
decision epoch. Its ten features are elapsed idle time, the current
position in the factory calendar, and the load history of the
production run that just ended.

Keeping them apart is not pedantry. Section~\ref{sec:explain} shows they
draw on almost disjoint information --- $11.6\%$ of the window model's
attribution is calendar against $82.4\%$ of the duration model's --- and
Section~\ref{sec:crossmachine} shows that one transfers between machines
and the other does not, for exactly that reason. Attributing the
economic result to the window model's $F_1$ would credit a model that
plays no part in producing it.

%  This file \inputs method_gmm_labeller.tex.
%  Bibliography: paper/references.bib
% =====================================================================

\section{Proposed Framework}
\label{sec:method}

\subsection{Overview}
\label{sec:method:overview}

The framework is a three-stage pipeline over a single input: the
per-phase voltage and current a standard retrofit meter already
records. Each stage is specified by what the next one needs.

\begin{enumerate}
\item \emph{State inference} (Section~\ref{sec:method:labelling}) maps
  the electrical record to a label per sample from
  $\{$OFF, STANDBY, WORKING, PEAK\_LOAD$\}$ without annotations, and is
  required to produce a sequence whose state durations are physically
  admissible --- not merely a per-sample partition --- because duration
  forecasting over a flickering label sequence is ill-posed.
\item \emph{Duration forecasting} (Sections~\ref{sec:forecasting}
  and~\ref{sec:decision:policies}) predicts how much longer the machine
  will remain non-productive, and is required to target the conditional
  \emph{mean} of that duration, because the decision rule below is
  affine in it.
\item \emph{Decision} (Section~\ref{sec:decision}) converts the forecast
  into a de-energise/keep-running action by minimising total cost
  subject to thermal and wear constraints, and returns the break-even
  duration as a derived quantity rather than consuming it as a
  parameter.
\end{enumerate}

The stages are separable in implementation and are evaluated separately
throughout, which is what makes the component ablation of
Section~\ref{sec:ablation} meaningful.

\subsection{Data}
\label{sec:method:data}

IMDELD~\cite{imdeld2018dataset,martins2018imdeld} publishes $1$\,Hz
three-phase measurements from eight machines of a Brazilian
poultry-feed pellet plant: two pelletizers, two milling machines, two
exhaust fans and two double-pole contactors. Six channels are used ---
active, reactive and apparent power, current, voltage and power factor
--- and no other information about the machines enters the models. The
plant's published calendar (weekday operation, a halt between 17:00 and
22:00 local time to avoid the peak tariff, and no weekend production)
is used \emph{only} as external validation evidence, never as a model
input to the state layer.

Coverage differs sharply by machine and is reported as measured rather
than as published. The pelletizer~I record spans $155$ days from
2017-10-30 to 2018-04-03 and contains $63.4$ days of data in $74$
contiguous segments; $31$ acquisition gaps account for the remainder,
the largest of them $30$ days. The two milling machines cover $12.2$ of
$43$ days in $1\,082$ and $580$ fragments. \textbf{Every annualised
quantity in this paper is scaled by covered time, not by span}, which
matters by a factor of $2.4$ on the five-month records and is the
reason the milling machines' figures are not annualised at all
(Section~\ref{sec:limitations:milling}).

% =====================================================================
%  method_gmm_labeller.tex  --  PHASE V CORRECTION 4 & 5
%
%  Drop-in description of the GMM-HMM labelling procedure, including
%  the accept/reject step that Phase V Task 12B showed is load-bearing
%  for reproducibility.  Insert into Section 3 (Method).
%
%  Numbers: all STANDBY-time totals are rounded to 2 significant
%  figures and stated as conditional on the physics battery passing,
%  per Phase V Task 11C (block-bootstrap CI ±15–25%) and
%  Task 12B (labelling CV 33.7% unconditional, 3.0% conditional).
% =====================================================================

\subsection{Machine-state inference}
\label{sec:method:labelling}

\subsubsection{Preprocessing}
\label{sec:method:preprocess}

Each machine's $1$\,Hz electrical record is preprocessed in three
deterministic steps, applied identically to every machine so that
differences in downstream results are attributable to the machines,
not to the preprocessing.

\begin{enumerate}
\item \emph{Spike masking.}  For each channel, rolling-MAD outliers
  (window $11$, threshold $5\times$MAD) are flagged and excluded
  from model fitting but retained in the record; removing them would
  destroy the row-count-to-seconds correspondence that episode
  durations depend on.

\item \emph{OFF-state denoising.}  Rows with active power below
  $5$\,W are smoothed with a rolling median (window $5$) to suppress
  sensor noise in the de-energised state without touching on-state
  dynamics.

\item \emph{Feature standardisation.}  Six electrical channels are
  log$_1$p-transformed (right-skewed channels: $P, Q, S, I$) and
  z-scored jointly before any model fitting.
\end{enumerate}

\subsubsection{Gaussian hidden Markov model}
\label{sec:method:ghmm}

Emission distributions are estimated by a Gaussian mixture model
(GMM, full covariance, $k = 4$ components) fitted on a uniform
random sample of $200{,}000$ rows.  A Gaussian HMM is then
warm-started from these emissions with $\texttt{params}=\texttt{``t''}$,
so only the transition matrix is learned from $64$ contiguous
two-hour blocks spread across the record; this decouples emission
shape (a property of the marginal distribution) from transition
dynamics (a property of the sequence) and avoids the local optima
that Baum--Welch finds when initialised randomly on this
dataset~\cite{fox2011sticky}.  Viterbi decoding is applied per
contiguous segment to prevent the decoder from asserting a transition
across an acquisition gap.  A minimum-dwell constraint of $10$\,s
is applied post-decode to remove label flicker without reshaping
the partition (Section~\ref{sec:results:task4}).

\subsubsection{Physics-validation accept/reject loop}
\label{sec:method:accept_reject}

\emph{This step is load-bearing for reproducibility and was not
explicit in the Phase I--IV description; it is stated here as a
result of Phase~V Task~12B.}

After decoding, a battery of ten physics tests (T1--T10:
power-factor separation, no-load current ratio, power conservation,
dwell plausibility, cross-machine replication, and factory-schedule
agreement) is run on a strided sample of the labelled record and
aggregated into a physics score $s \in [0, 1]$.  If $s < 0.90$,
the model is re-fitted with a different seed (incremented by a
fixed prime) and the battery is re-evaluated, up to three attempts.

The motivation is empirical.  Phase~V Task~12B fitted the labeller at
ten different seeds on identical preprocessed data from pelletizer~I
and found that three seeds (30\%) converge to a wrong cluster:

\begin{itemize}
\item seed~6: the ``STANDBY'' cluster absorbed a loaded running state
  (mean power $12{,}961$\,W instead of the correct $\sim3{,}750$\,W);
\item seeds~8 and~9: the ``STANDBY'' cluster absorbed the OFF state
  (mean power $0.8$--$1.1$\,W).
\end{itemize}

The physics battery scored those three seeds at $0.60$--$0.80$ and the
seven correct seeds at $1.00$, making the accept/reject decision exact
on this record.  Conditioning on a passing seed collapses the
STANDBY-hour coefficient of variation from $33.7\%$ (all ten seeds)
to $3.0\%$ (passing seeds), and the total reported throughout this
paper---$92 \pm 3$\,h over the 155-day record---is the mean of the
seven passing seeds.

\emph{All STANDBY-time totals in this paper are therefore stated as
conditional on the physics battery passing} (i.e., $s \ge 0.90$)
\emph{and are rounded to two significant figures}, because the
block-bootstrap sampling uncertainty (Phase~V Task~11C, one-week
blocks) is $\pm 15$--$25\%$ and the conditional labelling
uncertainty is $\pm 3\%$.  Three significant figures are not
supported by the data.

The accept/reject loop is implemented in \texttt{src/labelling.py}
(\texttt{label\_record}, \texttt{physics\_score\_threshold}=0.90,
\texttt{max\_refit\_attempts}=3) and its output is stored in the
model-info dictionary (\texttt{physics\_validation\_attempts},
\texttt{physics\_validation\_seed}, \texttt{physics\_score}) for
audit and reproduction.


\subsection{Two forecasting problems, not one}
\label{sec:method:forecasting}

The framework contains two forecasters. They are frequently conflated,
they consume different features, and no result in this paper reports
them in the same column.

\emph{The window state forecaster} takes a $600$\,s lookback of the
electrical record and predicts the state at each step of a $300$\,s
horizon. It is a sequence problem, it is what the baseline comparison of
Section~\ref{sec:forecasting} evaluates, and it is scored on windows.
Its design representation summarises each channel over the lookback ---
final observed value, sub-block means and standard deviations --- giving
$217$ columns for the tree, or the raw sequence for the neural
architectures.

\emph{The episode-duration forecaster} takes the state of an idle
episode at a decision epoch and predicts the \emph{remaining} idle time
$R$ in seconds. It is a regression problem, it is what the decision
layer of Section~\ref{sec:decision} consumes, and it is scored per
decision epoch. Its ten features are elapsed idle time, the current
position in the factory calendar, and the load history of the
production run that just ended.

Keeping them apart is not pedantry. Section~\ref{sec:explain} shows they
draw on almost disjoint information --- $11.6\%$ of the window model's
attribution is calendar against $82.4\%$ of the duration model's --- and
Section~\ref{sec:crossmachine} shows that one transfers between machines
and the other does not, for exactly that reason. Attributing the
economic result to the window model's $F_1$ would credit a model that
plays no part in producing it.
